\documentclass[11pt,a4paper]{article}

% ── Codificación y fuentes ────────────────────────────────────────────────────
\usepackage[utf8]{inputenc}
\usepackage[T1]{fontenc}
\usepackage{lmodern}
\usepackage[spanish]{babel}

% ── Geometría y espaciado ─────────────────────────────────────────────────────
\usepackage[top=2.5cm, bottom=2.5cm, left=2.8cm, right=2.8cm]{geometry}
\usepackage{setspace}
\setstretch{1.15}
\usepackage{parskip}

% ── Tablas ────────────────────────────────────────────────────────────────────
\usepackage{booktabs}
\usepackage{tabularx}
\usepackage{array}
\usepackage{longtable}
\usepackage{enumitem}

% ── Colores y cajas ───────────────────────────────────────────────────────────
\usepackage[dvipsnames]{xcolor}
\usepackage{tcolorbox}
\tcbuselibrary{skins, breakable}

% ── Cabeceras y pies de página ────────────────────────────────────────────────
\usepackage{fancyhdr}
\usepackage{lastpage}
\pagestyle{fancy}
\fancyhf{}
\fancyhead[L]{\small\textbf{Informe de Evaluación} --- Física para Bioanalistas}
\fancyhead[R]{\small Daniela Escribano}
\fancyfoot[C]{\small Página \thepage\ de \pageref{LastPage}}
\renewcommand{\headrulewidth}{0.4pt}

% ── Hiperenlaces ──────────────────────────────────────────────────────────────
\usepackage[colorlinks=true, linkcolor=NavyBlue, urlcolor=NavyBlue]{hyperref}

% ── Paleta de colores personalizados ─────────────────────────────────────────
\definecolor{desempeno_excelente}{HTML}{1A6B2F}
\definecolor{desempeno_bueno}{HTML}{2E5A9C}
\definecolor{desempeno_suficiente}{HTML}{8B6914}
\definecolor{desempeno_deficiente}{HTML}{C0392B}
\definecolor{desempeno_insuficiente}{HTML}{7B241C}
\definecolor{grisFondo}{HTML}{F5F5F5}
\definecolor{azulTitulo}{HTML}{1B2A6B}
\definecolor{verdeFortaleza}{HTML}{1A6B2F}
\definecolor{naranjaMejora}{HTML}{A04000}

% ── Estilos de cajas ─────────────────────────────────────────────────────────
\tcbset{
  cajaTitulo/.style={
    enhanced, breakable,
    colback=azulTitulo!8, colframe=azulTitulo,
    fonttitle=\bfseries\large, coltitle=white,
    attach boxed title to top left={yshift=-2mm, xshift=4mm},
    boxed title style={colback=azulTitulo},
    top=6pt, bottom=6pt, left=8pt, right=8pt,
  },
  cajaFortaleza/.style={
    enhanced, breakable,
    colback=verdeFortaleza!6, colframe=verdeFortaleza!60,
    fonttitle=\bfseries, coltitle=verdeFortaleza,
    top=4pt, bottom=4pt, left=8pt, right=8pt,
  },
  cajaMejora/.style={
    enhanced, breakable,
    colback=naranjaMejora!6, colframe=naranjaMejora!60,
    fonttitle=\bfseries, coltitle=naranjaMejora,
    top=4pt, bottom=4pt, left=8pt, right=8pt,
  },
  cajaRecomendacion/.style={
    enhanced, breakable,
    colback=grisFondo, colframe=gray!50,
    fonttitle=\bfseries, coltitle=black,
    top=4pt, bottom=4pt, left=8pt, right=8pt,
  },
}

% ─────────────────────────────────────────────────────────────────────────────
\begin{document}

% ══ PORTADA ══════════════════════════════════════════════════════════════════
\begin{center}
  {\Large\bfseries\color{azulTitulo} Universidad de Oriente/Núcleo Sucre --- Física para Bioanalistas}\\[4pt]
  {\large Informe de Evaluación Preliminar}\\[12pt]
  \rule{\linewidth}{1.2pt}\\[6pt]
  {\LARGE\bfseries Daniela Escribano}\\[4pt]
  \rule{\linewidth}{0.6pt}
\end{center}

% ══ FICHA TÉCNICA ════════════════════════════════════════════════════════════
\vspace{4pt}
\begin{tcolorbox}[cajaTitulo, title=Datos del informe]
\begin{tabular}{@{}llll@{}}
  \textbf{Estudiante:}  & Daniela Escribano  &
  \textbf{Fecha:}       & 2026-08-08 \\[4pt]
  \textbf{Archivo PDF:} & \texttt{Daniela\_Escribano.pdf}  &
  \textbf{Páginas:}     & 7 \\
\end{tabular}
\end{tcolorbox}

% ══ RESULTADO GLOBAL ═════════════════════════════════════════════════════════
\vspace{6pt}
\begin{center}
  \begin{tcolorbox}[
    enhanced,
    colback=white, colframe=desempeno_suficiente,
    width=0.62\linewidth,
    boxrule=2pt,
    halign=center, valign=center,
    top=8pt, bottom=8pt,
  ]
    \centering
    {\large\bfseries Puntaje sugerido}\\[6pt]
    {\Huge\bfseries\color{desempeno_suficiente} 6.8/10}\\[6pt]
    {\large Nivel de desempeño: \textbf{\color{desempeno_suficiente} Suficiente}}
  \end{tcolorbox}
\end{center}

% ══ RESUMEN GENERAL ══════════════════════════════════════════════════════════
\section*{Resumen general}
El estudiante demuestra un buen dominio de los cálculos numéricos y la aplicación de fórmulas en los problemas que abordó. Sin embargo, el desempeño se ve afectado por la falta de respuesta en varias preguntas cualitativas (4a, 4b, 5a, 5b) y por explicaciones conceptuales incompletas en otras. Las partes respondidas muestran una comprensión adecuada de los principios físicos, pero la evaluación general indica la necesidad de fortalecer la argumentación y el análisis cualitativo de los fenómenos biológicos.

% ══ FORTALEZAS Y ÁREAS DE MEJORA ════════════════════════════════════════════
\vspace{4pt}
\begin{minipage}[t]{0.48\linewidth}
  \begin{tcolorbox}[cajaFortaleza, title={Fortalezas}]
\begin{itemize}[leftmargin=1.5em, itemsep=2pt]
  \item Excelente habilidad para realizar cálculos numéricos y aplicar fórmulas correctamente.
  \item Manejo preciso de unidades y conversiones en los problemas resueltos.
  \item Claridad y organización en la presentación de los datos y el desarrollo de los cálculos.
  \item Buena comprensión conceptual en las preguntas 1a, 1b, 2b y 3a, donde se justifican los fenómenos físicos.
\end{itemize}
  \end{tcolorbox}
\end{minipage}
\hfill
\begin{minipage}[t]{0.48\linewidth}
  \begin{tcolorbox}[cajaMejora, title={Areas de mejora}]
\begin{itemize}[leftmargin=1.5em, itemsep=2pt]
  \item Completitud de las respuestas: varias sub-preguntas cualitativas fueron dejadas en blanco.
  \item Profundización en la argumentación y análisis físico de los fenómenos, especialmente en las consecuencias biológicas y clínicas.
  \item Mejorar la capacidad de articular explicaciones cualitativas detalladas, no solo las fórmulas o los datos.
\end{itemize}
  \end{tcolorbox}
\end{minipage}

% ══ OBSERVACIONES POR PREGUNTA ═══════════════════════════════════════════════
\section*{Observaciones por pregunta}

\begin{longtable}{>{\bfseries}p{2.8cm} p{1.6cm} p{7.0cm} p{2.8cm}}
  \toprule
  \textbf{Pregunta} & \textbf{Puntaje} & \textbf{Evaluación} & \textbf{Errores detectados} \\
  \midrule
  \endhead
  \bottomrule
  \endfoot
    \textbf{Pregunta 1: Termorregulación y Eficiencia de la Sudoración} & 2.0/2.0 & El estudiante responde de manera excelente a todas las partes de esta pregunta. La justificación de la diferencia entre los pacientes A y B es clara y conceptualmente correcta, destacando el rol del calor latente de vaporización. La explicación del efecto de la humedad relativa al 100\% es precisa, y el cálculo del calor disipado es impecable, incluyendo la conversión a kilocalorías. & \textit{---} \\
    \midrule
    \textbf{Pregunta 2: Mecánica Respiratoria y Ley de Laplace} & 1.7/2.0 & En la parte (a), el estudiante explica correctamente la Ley de Laplace y la relación inversa entre presión y radio cuando la tensión superficial es constante. Sin embargo, omite la consecuencia crucial de este fenómeno: el flujo de aire de alvéolos pequeños a grandes y el colapso de los alvéolos pequeños. La parte (b) sobre el surfactante es muy buena, explicando su función y mecanismo de acción. El cálculo en la parte (c) es correcto. & \textit{Error conceptual: Omisión de la consecuencia del colapso alveolar en la parte (a) cuando la tensión superficial es constante.} \\
    \midrule
    \textbf{Pregunta 3: Optimización en Hemodiálisis y Primera Ley de Fick} & 1.9/2.0 & La parte (a) está muy bien resuelta, aplicando correctamente la Ley de Fick para determinar el aumento de eficiencia. La parte (b) explica adecuadamente el riesgo mecánico de reducir el espesor de la membrana, aunque podría haber detallado más las consecuencias biológicas (hemólisis, infección). El cálculo en la parte (c) es correcto y bien presentado. & \textit{Explicación incompleta: En la parte (b), no se mencionan explícitamente las consecuencias clínicas como hemólisis o infección por la ruptura de la membrana.} \\
    \midrule
    \textbf{Pregunta 4: Errores de Infusión Intravenosa y Ósmosis} & 0.6/2.0 & Las partes (a) y (b) de esta pregunta fueron dejadas en blanco, lo que representa una omisión significativa de contenido conceptual clave sobre ósmosis y sus efectos en los eritrocitos. La parte (c), el cálculo de la presión osmótica, es correcta en su fórmula, datos y resultado. & \textit{Pregunta en blanco: Las partes (a) y (b) no fueron respondidas.} \\
    \midrule
    \textbf{Pregunta 5: Capilaridad y Toma de Muestras Sanguíneas} & 0.6/2.0 & Las partes (a) y (b) de esta pregunta fueron dejadas en blanco, lo que indica una falta de conocimiento o de tiempo para abordar las explicaciones cualitativas sobre capilaridad. El cálculo de la altura de ascenso capilar en la parte (c) es correcto, con la fórmula, datos y resultado bien presentados. & \textit{Pregunta en blanco: Las partes (a) y (b) no fueron respondidas.} \\
    \midrule
\end{longtable}

% ══ ERRORES DETECTADOS ═══════════════════════════════════════════════════════
\section*{Análisis de errores}

\subsection*{Errores conceptuales}
\begin{itemize}[leftmargin=1.5em, itemsep=2pt]
  \item En la Pregunta 2a, el estudiante no describe la consecuencia directa de la Ley de Laplace en los alvéolos (flujo de aire de pequeños a grandes y colapso de los pequeños) cuando la tensión superficial es constante.
  \item Omisión de respuestas en las preguntas 4a, 4b, 5a y 5b, lo que sugiere lagunas conceptuales en ósmosis y capilaridad.
\end{itemize}

\subsection*{Errores procedimentales}
\textit{Ninguno identificado.}

% ══ RECOMENDACIONES AL ESTUDIANTE ════════════════════════════════════════════
\section*{Retroalimentación para el estudiante}
\begin{tcolorbox}[cajaRecomendacion, title={Dirigido a: Daniela Escribano}]
Estimado estudiante, has demostrado una sólida capacidad para resolver problemas numéricos y aplicar fórmulas de física. Sin embargo, es crucial que dediques más tiempo a comprender y articular las explicaciones cualitativas y las implicaciones biológicas de los principios físicos. Asegúrate de responder todas las partes de cada pregunta, incluso si son explicativas. Practica la redacción de argumentos completos que conecten la física con los fenómenos clínicos, como los efectos de la ósmosis en las células sanguíneas o las consecuencias de la Ley de Laplace en la respiración. Revisar los conceptos teóricos y practicar la argumentación teórica te ayudará a mejorar significativamente tu desempeño.
\end{tcolorbox}

% ══ NOTA PARA EL PROFESOR ════════════════════════════════════════════════════
\section*{Nota para el profesor}
\begin{tcolorbox}[
  enhanced, breakable,
  colback=yellow!8, colframe=yellow!60!black,
  fonttitle=\bfseries, coltitle=black,
  title={Observaciones para revisión manual},
  top=4pt, bottom=4pt, left=8pt, right=8pt,
]
El estudiante presenta un buen nivel en la resolución de problemas cuantitativos, aplicando correctamente las fórmulas y manejando las unidades. Sin embargo, hay una debilidad notable en las preguntas de índole cualitativa y explicativa, con varias sub-preguntas dejadas en blanco (4a, 4b, 5a, 5b) y una explicación incompleta en 2a. Esto sugiere que, si bien el estudiante puede aplicar las matemáticas, necesita reforzar la comprensión conceptual profunda y la capacidad de argumentación física aplicada a contextos biológicos, tal como se enfatiza en la rúbrica. El puntaje refleja un buen desempeño en lo que se intentó, pero penaliza fuertemente las omisiones.
\end{tcolorbox}

% ══ PIE DE INFORME ════════════════════════════════════════════════════════════
\vspace{12pt}
\begin{center}
  \footnotesize\color{gray}
  Informe generado automáticamente por el sistema de evaluación con IA (gemini-2.5-flash) y soporte LaTex. \\
  Este documento es una evaluación \textbf{preliminar} para ser revisado por el profesor antes de ser entregado al 
  estudiante. \\
  Generado el 2026-08-08 \textbullet\ BIOANALISIS Sistema Académico de Evaluación.
  \end{center}

\end{document}
